\subsection{Quantitative Evaluation of Segmentation Performance}

\begin{table}[t]
\centering
\caption{Comparison of CATMIL with standard losses on the MSLesSeg test
set, without post-processing. Values are the mean $\pm$ sample standard
deviation over the five fold models. Best values are shown in
\textbf{bold}; second-best values are \underline{underlined}.}
\label{tab:loss_comparison_small_lesion}
\resizebox{\textwidth}{!}{%
\begin{tabular}{lcccc}
\toprule
\textbf{Metric}
& \textbf{CATMIL (ours)}
& \textbf{DiceCE}
& \textbf{Tversky}
& \textbf{FocalTversky} \\
\midrule
Small Lesion Recall $\uparrow$
& \textbf{0.8730 $\pm$ 0.01}
& 0.7956 $\pm$ 0.03
& \underline{0.8336 $\pm$ 0.02}
& 0.8313 $\pm$ 0.02 \\
FN Count $\downarrow$
& \textbf{2.5667 $\pm$ 0.15}
& 4.9500 $\pm$ 0.40
& 3.7667 $\pm$ 0.23
& \underline{3.7000 $\pm$ 0.24} \\
FN Volume Fraction $\downarrow$
& \textbf{0.0214 $\pm$ 0.003}
& 0.0341 $\pm$ 0.01
& 0.0257 $\pm$ 0.01
& \underline{0.0250 $\pm$ 0.005} \\
\midrule
DSC $\uparrow$
& \textbf{0.7834 $\pm$ 0.01}
& 0.7796 $\pm$ 0.01
& 0.7706 $\pm$ 0.01
& \underline{0.7802 $\pm$ 0.01} \\
HD95 (mm) $\downarrow$
& \textbf{7.9817 $\pm$ 1.84}
& 9.0372 $\pm$ 1.34
& 10.2408 $\pm$ 0.15
& \underline{8.2060 $\pm$ 1.93} \\
Lesion F1 $\uparrow$
& 0.7571 $\pm$ 0.02
& \underline{0.8433 $\pm$ 0.01}
& 0.8402 $\pm$ 0.01
& \textbf{0.8455 $\pm$ 0.01} \\
FP Volume (mm$^3$) $\downarrow$
& \textbf{1537 $\pm$ 245}
& \underline{1819 $\pm$ 654}
& 2621 $\pm$ 559
& 2282 $\pm$ 710 \\
\bottomrule
\end{tabular}%
}
\end{table}

\Cref{tab:loss_comparison_small_lesion} summarizes the quantitative
performance of the evaluated loss functions on the independent MSLesSeg
test set, and \cref{tab:paired_differences} gives the paired differences
between CATMIL and each baseline with patient-level 95\% CIs
(\cref{sec:stats}). CATMIL has the best mean value in six of the seven
metrics, but the paired analysis shows that only some of these differences
are supported. The detection metrics improve consistently across patients.
DSC and HD95 are comparable to those of the baselines, and lesion-wise F1
is consistently lower.

\begin{table}[t]
\centering
\caption{Paired differences between CATMIL and each baseline (CATMIL minus
baseline), with 95\% patient-level bootstrap CIs and, in parentheses, the
number of test patients in which CATMIL is better. Differences whose CI
excludes zero are shown in \textbf{bold}. PP5: after removing predicted
components smaller than 5 voxels. With six patients, the smallest possible
Holm-corrected Wilcoxon $p$-value is 0.094 (\cref{sec:stats}).}
\label{tab:paired_differences}
\resizebox{\textwidth}{!}{%
\begin{tabular}{llll}
\toprule
\textbf{Metric} & \textbf{vs.\ DiceCE} & \textbf{vs.\ Tversky} & \textbf{vs.\ FocalTversky} \\
\midrule
\multicolumn{4}{l}{\emph{MSLesSeg}} \\
Small Lesion Recall $\uparrow$ & \textbf{+0.077 [+0.030, +0.157] (6/6)} & \textbf{+0.039 [+0.004, +0.093] (6/6)} & \textbf{+0.042 [+0.005, +0.101] (6/6)} \\
FN Count $\downarrow$ & \textbf{$-$2.38 [$-$5.29, $-$0.84] (6/6)} & \textbf{$-$1.20 [$-$2.98, $-$0.27] (6/6)} & \textbf{$-$1.13 [$-$2.71, $-$0.27] (5/6)} \\
DSC $\uparrow$ & +0.004 [$-$0.010, +0.031] (1/6) & \textbf{+0.013 [+0.001, +0.034] (5/6)} & +0.003 [$-$0.015, +0.023] (5/6) \\
HD95 (mm) $\downarrow$ & $-$1.06 [$-$4.67, +0.73] (4/6) & $-$2.26 [$-$7.35, +0.11] (4/6) & $-$0.22 [$-$2.92, +2.72] (4/6) \\
Lesion F1 $\uparrow$ & \textbf{$-$0.086 [$-$0.138, $-$0.043] (0/6)} & \textbf{$-$0.083 [$-$0.143, $-$0.038] (0/6)} & \textbf{$-$0.088 [$-$0.151, $-$0.039] (0/6)} \\
FP Volume (mm$^3$) $\downarrow$ & $-$282 [$-$1582, +271] (1/6) & \textbf{$-$1084 [$-$2675, $-$302] (6/6)} & \textbf{$-$746 [$-$2189, $-$182] (6/6)} \\
\midrule
\multicolumn{4}{l}{\emph{MSLesSeg, PP5}} \\
Small Lesion Recall $\uparrow$ & \textbf{+0.070 [+0.023, +0.151] (6/6)} & \textbf{+0.033 [+0.001, +0.079] (6/6)} & \textbf{+0.036 [+0.000, +0.093] (6/6)} \\
FN Count $\downarrow$ & \textbf{$-$2.10 [$-$4.80, $-$0.66] (6/6)} & \textbf{$-$0.93 [$-$2.38, $-$0.13] (6/6)} & \textbf{$-$0.95 [$-$2.40, $-$0.14] (6/6)} \\
DSC $\uparrow$ & +0.004 [$-$0.009, +0.032] (2/6) & \textbf{+0.013 [+0.001, +0.035] (5/6)} & +0.004 [$-$0.014, +0.024] (5/6) \\
HD95 (mm) $\downarrow$ & $-$0.99 [$-$4.65, +0.76] (4/6) & $-$2.26 [$-$7.25, +0.10] (5/6) & $-$0.13 [$-$2.80, +2.76] (4/6) \\
Lesion F1 $\uparrow$ & $-$0.017 [$-$0.040, +0.001] (2/6) & $-$0.012 [$-$0.038, +0.006] (2/6) & $-$0.017 [$-$0.043, +0.001] (2/6) \\
FP Volume (mm$^3$) $\downarrow$ & $-$301 [$-$1610, +244] (1/6) & \textbf{$-$1103 [$-$2706, $-$310] (6/6)} & \textbf{$-$765 [$-$2211, $-$197] (6/6)} \\
\midrule
\multicolumn{4}{l}{\emph{3D-MR-MS}} \\
Small Lesion Recall $\uparrow$ & \textbf{+0.081 [+0.034, +0.127] (5/6)} & \textbf{+0.054 [+0.019, +0.086] (5/6)} & \textbf{+0.055 [+0.021, +0.087] (5/6)} \\
FN Count $\downarrow$ & \textbf{$-$5.13 [$-$12.40, $-$0.50] (5/6)} & \textbf{$-$4.43 [$-$11.23, $-$0.27] (5/6)} & \textbf{$-$4.43 [$-$11.47, $-$0.23] (5/6)} \\
DSC $\uparrow$ & $-$0.012 [$-$0.032, +0.010] (0/6) & $-$0.008 [$-$0.030, +0.013] (3/6) & $-$0.016 [$-$0.035, +0.001] (1/6) \\
HD95 (mm) $\downarrow$ & $-$0.38 [$-$4.14, +2.07] (3/6) & $-$1.83 [$-$8.51, +1.50] (2/6) & $-$0.72 [$-$6.07, +2.92] (3/6) \\
Lesion F1 $\uparrow$ & \textbf{$-$0.197 [$-$0.258, $-$0.142] (0/6)} & \textbf{$-$0.168 [$-$0.242, $-$0.106] (0/6)} & \textbf{$-$0.166 [$-$0.232, $-$0.108] (0/6)} \\
FP Volume (mm$^3$) $\downarrow$ & $-$2 [$-$110, +110] (3/6) & \textbf{$-$761 [$-$1587, $-$118] (6/6)} & \textbf{$-$726 [$-$1550, $-$92] (6/6)} \\
\bottomrule
\end{tabular}%
}
\end{table}

\textbf{CATMIL is more sensitive to small lesions than all three standard
loss functions.} It achieves a Small Lesion Recall of $0.8730 \pm 0.01$,
compared with $0.7956 \pm 0.03$ for DiceCE, $0.8336 \pm 0.02$ for Tversky,
and $0.8313 \pm 0.02$ for FocalTversky. CATMIL therefore improves
small-lesion recall by 7.74 percentage points over DiceCE and by 3.94
percentage points over the next-best baseline, Tversky. In practical terms,
CATMIL detects approximately 87\% of small lesions, whereas the competing
methods detect approximately 80--83\%. The gain is present in every test
patient against every baseline, and its 95\% CI excludes zero in each
comparison (\cref{tab:paired_differences}): +0.077 [+0.030, +0.157] against
DiceCE and +0.039 [+0.004, +0.093] against Tversky. The gain over DiceCE
does not depend on the lenient any-overlap detection criterion: when at
least half of a lesion's voxels must be covered, it remains 8 percentage
points on MSLesSeg and 10 on 3D-MR-MS, with CIs that exclude zero
(\cref{appendix:strict-detection}).

\textbf{CATMIL misses fewer lesions and less lesion tissue.} CATMIL
produces an average FN Count of $2.57 \pm 0.15$, compared with
$4.95 \pm 0.40$ for DiceCE, corresponding to approximately 48\% fewer
completely missed lesions. It also achieves the lowest FN Volume Fraction
($0.0214 \pm 0.003$), compared with $0.0341 \pm 0.01$ for DiceCE. Thus, the
remaining false-negative regions are fewer in number and occupy less total
volume. The reduction in missed lesions is present in all six patients
against DiceCE and Tversky and in five of six against FocalTversky, and the
small standard deviations show that it does not depend on one favorable
fold.

\textbf{CATMIL improves detection without increasing false-alarm volume.}
It produces the lowest mean FP Volume ($1537 \pm 245$~mm$^3$), compared
with $1819 \pm 654$~mm$^3$ for DiceCE, $2621 \pm 559$~mm$^3$ for Tversky,
and $2282 \pm 710$~mm$^3$ for FocalTversky. The reduction relative to the
Tversky-based losses holds in all six patients, whereas the difference from
DiceCE is not supported ($-282$~mm$^3$, 95\% CI [$-1582$, +271]). CATMIL
therefore does not improve recall by predicting a larger volume of lesion
tissue.

\textbf{CATMIL preserves overlap and boundary accuracy.} Its mean DSC
($0.7834 \pm 0.01$) and HD95 ($7.98 \pm 1.84$~mm) are the best of the four
losses, but the differences are small relative to the variation between
patients. The 95\% CIs include zero for every comparison except DSC against
Tversky, and the DSC of CATMIL is higher than that of DiceCE in only one of
the six patients (\cref{tab:paired_differences}). The additional lesion
detections are therefore not obtained at the cost of poorer overlap or less
accurate boundaries, but CATMIL does not improve these metrics either.

\textbf{CATMIL's main limitation is its lesion-wise precision--recall
balance rather than its ability to detect or segment lesions.} CATMIL
achieves a Lesion F1 score of $0.7571 \pm 0.02$, which is lower than
$0.8433 \pm 0.01$ for DiceCE, $0.8402 \pm 0.01$ for Tversky, and
$0.8455 \pm 0.01$ for FocalTversky, that is, 8.84 percentage points below
the best result, and it is lower in every test patient. Because Lesion F1
is affected by the number of
false-positive lesion components, multiple small disconnected predictions
can reduce this metric substantially even when their total volume is
limited. The combination of high small-lesion recall, low FP Volume, and
comparatively low Lesion F1 therefore suggests that CATMIL produces
numerous small false-positive components rather than large over-segmented
regions. This hypothesis is examined in \cref{sec:fp_pp}.

\subsection{Analysis of Missed Lesion Errors Across Sizes}

CATMIL leaves the fewest lesions completely undetected. In relative terms,
CATMIL misses 9.2\% of the reference lesions, compared with 12.0\% for
Tversky and FocalTversky and 15.3\% for DiceCE. This subsection examines
which lesions account for this difference and whether the reduction holds
consistently across test cases.

\Cref{fig:recall_vs_size_curve} shows lesion recall as a function of lesion
size. The CATMIL curve lies above all three baselines over the whole
small-lesion range, and the gap is widest where the baselines fail most.
For lesions of 5--30~mm$^3$, CATMIL detects 87\% of lesions, compared with
57\% for DiceCE and 72\% for Tversky and FocalTversky. Over the bins of the
size analysis, CATMIL raises recall from 0.12--0.29 to 0.45 for lesions of
at most 10~mm$^3$ and from 0.74--0.82 to 0.89 for lesions of
11--50~mm$^3$. Above about 100~mm$^3$ all curves converge near 1, so the
gain on small lesions is not paid for with large ones. CATMIL therefore
shifts the size at which lesions become reliably detectable toward smaller
lesions: it reaches a recall of 0.8 at about 13~mm$^3$, whereas Tversky and
FocalTversky reach it only at about 25~mm$^3$ and DiceCE at about
40~mm$^3$. The smallest lesions, of at most 5~mm$^3$, remain difficult for
every loss, and the size bias of voxel-wise training is reduced but not
removed.

\begin{figure}[t]
    \centering
    \includegraphics[width=0.85\textwidth]{figures/recall_vs_size_curve.pdf}
    \caption{Lesion recall against lesion size on the MSLesSeg test set.
    Upper panel: fraction of reference lesions detected in each size
    interval (logarithmic axis); lines are the mean over the five fold
    models and shaded bands span the fold minimum to maximum. Lower panel:
    number of reference lesions per interval. Dotted vertical lines mark
    10, 50, and 200~mm$^3$. A lesion counts as detected if any predicted
    voxel overlaps it.}
    \label{fig:recall_vs_size_curve}
\end{figure}

A higher recall could in principle hide a trade, in which CATMIL gains some
lesions but loses others that a baseline detects.
\Cref{fig:lesion_fate} tests this by pairing every reference lesion between
CATMIL and each baseline trained on the same fold. The comparison is
strongly one-sided. Against DiceCE, CATMIL detects 30.6 lesions per fold
model that DiceCE misses, while DiceCE detects only 2.0 that CATMIL misses.
Against Tversky and FocalTversky the ratios are 17.2 to 2.8 and 17.0 to
3.4. About two thirds of the lesions recovered by CATMIL are at most
50~mm$^3$, and no lesion above 200~mm$^3$ is detected by a baseline but
missed by CATMIL. The recall gain is therefore a net gain in detected
lesions, not a redistribution of errors between lesion sizes.

\begin{figure}[t]
    \centering
    \includegraphics[width=\textwidth]{figures/lesion_fate.pdf}
    \caption{Lesions detected by only one of two losses on the MSLesSeg
    test set, by lesion size. Each panel pairs every reference lesion
    between CATMIL and one baseline trained on the same fold. Bars to the
    right: lesions detected by CATMIL but not by the baseline; bars to the
    left: lesions detected by the baseline but not by CATMIL. Grey text:
    lesions detected by both or by neither. Values are means over the five
    fold models.}
    \label{fig:lesion_fate}
\end{figure}

The recovered lesions raise the question of why the baselines miss them:
the baseline output may contain a weak response that falls just below the
decision threshold, or no response at all. \Cref{fig:max_prob_by_size}
distinguishes the two cases by assigning the maximum foreground
probability inside every reference lesion to one of three levels: below
0.1 (no response), between 0.1 and 0.5 (a weak response below the
threshold), or at least 0.5 (detected).

For the smallest lesions, the baselines produce no response to most of
them. Under DiceCE, 82\% of lesions of at most 10~mm$^3$ have a maximum
probability below 0.1, and under the Tversky-based losses 65--68\%. CATMIL
reduces this fraction to 47\%, and for lesions of 11--50~mm$^3$ from
17--21\% to 8\%. The intermediate level is small for every loss (at most
8\%), so the lesions missed by the baselines cannot be recovered by
lowering the threshold. CATMIL therefore does not merely push weak
responses over the threshold: it makes the network respond to small
lesions that a voxel-wise objective leaves invisible. This is the effect
that the two terms of CATMIL are designed to produce: the CAT term gives
small and large lesions comparable total weight, and the MIL term
penalizes any lesion in which no voxel receives a high probability.

\begin{figure}[t]
    \centering
    \includegraphics[width=0.75\textwidth]{figures/max_prob_by_size.pdf}
    \caption{Maximum foreground probability inside each reference lesion on
    the MSLesSeg test set, by lesion size. Bars give the fraction of
    lesions whose maximum probability is below 0.1 (no response), between
    0.1 and 0.5 (weak response below the threshold), or at least 0.5
    (detected). Each lesion contributes once per fold model; $n$ is the
    number of reference lesions in the bin.}
    \label{fig:max_prob_by_size}
\end{figure}

Averages over the test set can hide whether an improvement is shared
across patients or produced by a few favorable cases.
\Cref{fig:fn_per_case_paired} compares CATMIL with each baseline case by
case. CATMIL does not trade errors between cases. Against DiceCE, no test
case becomes worse. Against the Tversky-based losses, the few cases in
which CATMIL misses more lesions differ by less than half a lesion, well
within the variation between trained models. The gain is instead
concentrated in the high-burden cases P47\_T1, P49\_T1, and P13\_T2, which
contain many tiny and small lesions (\cref{tab:test_case_lesion_profile}).
About 78\% of the lesions that CATMIL recovers relative to DiceCE in these
cases are at most 50 voxels in size. The same cases also explain why this
effect is barely visible in DSC. In P47\_T1, CATMIL recovers about 11
lesions on average relative to DiceCE, while the case-level DSC remains
unchanged at 0.85, because the recovered lesions are too small to change a
volume-overlap score. Because the 12 test cases come from six patients,
with several follow-up scans of the same patient, the case-level pattern
is complemented by the patient-level analysis in
\cref{tab:paired_differences}, in which CATMIL misses fewer lesions than
DiceCE in all six patients.

\begin{figure}[t]
    \centering
    \includegraphics[width=\textwidth]{figures/fn_per_case_paired.pdf}
    \caption{Missed lesions per test case for CATMIL compared with DiceCE,
    Tversky, and FocalTversky. Each point is one test case, placed at the
    mean FN lesion count over the five fold models; error bars span the
    fold minimum to maximum. Points below the dashed diagonal (shaded
    region) are cases in which CATMIL misses fewer lesions than the
    baseline. The text in each panel gives the number of cases in which
    CATMIL misses fewer, the same number, or more lesions, and the mean
    number of missed lesions per case for the baseline and for CATMIL.}
    \label{fig:fn_per_case_paired}
\end{figure}

\subsection{Case-Based Qualitative Analysis of Lesion Segmentation}

The size analysis shows how many lesions CATMIL recovers, but not which
lesions they are or why the baselines lose them. This subsection therefore
examines individual lesions rather than whole slices: at full-slice scale,
the lesions that separate the losses occupy only a few pixels.
\Cref{fig:qual_lesion_gallery} shows each lesion in a 24~mm crop centred on
it, and \cref{fig:qual_probability_maps} shows the lesion probabilities
that each loss assigns to the same regions.

The lesions were selected by a fixed rule applied to the per-lesion
detection results of all five fold models, not by visual inspection. A
\emph{recovered} lesion has at most 50 voxels, is detected by at least four
of the five CATMIL models, and is detected by at most one of the five
models of each baseline. Four such lesions are shown, from cases P47\_T1
and P26\_T1, both of which contain many small lesions. A lesion
\emph{missed by all} is detected by none of the twenty models. The
\emph{CATMIL false positive} is a predicted component produced by all five
CATMIL models and by none of the fifteen baseline models. Each lesion is
characterized by its size, its axial extent, its local FLAIR contrast
(computed as in \cref{appendix:test-case-lesion-profile}), its distance to
the nearest other lesion, and its position relative to the cortex.
Contrast values are also expressed as percentiles among the 84 test lesions
of 5--50 voxels.

\begin{figure}[!htbp]
    \centering
    \includegraphics[width=\textwidth]{figures/qual_lesion_gallery.pdf}
    \caption{Lesion-centered comparison of the four losses on the MSLesSeg
    test set. Each row shows one lesion (or one CATMIL false-positive
    component), selected by the rule described in the text. From left to
    right: the full axial slice with the crop location (cyan box); a 24~mm
    FLAIR crop with the target marked by a cyan arrow; the reference
    contour (white); and the predictions of the fold-0 model of each loss.
    Green indicates correctly segmented lesion voxels, red indicates missed
    lesion voxels (false negatives), and yellow indicates false-positive
    voxels. The number at the bottom left of each prediction panel is the
    number of the five fold models that detect the lesion; $p$ at the
    bottom right is the maximum lesion probability inside the lesion,
    averaged over the five fold models. The models receive T1-w, T2-w, and
    FLAIR images; only FLAIR is shown.}
    \label{fig:qual_lesion_gallery}
\end{figure}

\begin{figure}[!htbp]
    \centering
    \includegraphics[width=0.8\textwidth]{figures/qual_probability_maps.pdf}
    \caption{Lesion probability maps for the lesions in
    \cref{fig:qual_lesion_gallery}, averaged over the five fold models of
    each loss. The cyan contour marks the reference lesion and the yellow
    contour marks the CATMIL false-positive component. Values of $p$ are as
    in \cref{fig:qual_lesion_gallery}.}
    \label{fig:qual_probability_maps}
\end{figure}

\textbf{The baselines do not narrowly miss the recovered lesions; they do
not respond to them.} For three of the four recovered lesions, the maximum
DiceCE probability inside the lesion is at most 0.01, and for the 13- and
21-voxel lesions in P47\_T1, the Tversky and FocalTversky models also
assign a probability of 0.00 (\cref{fig:qual_probability_maps}). A lower
decision threshold would therefore not recover these lesions, because the
omission is already present in the network output rather than being
introduced by the decision rule. CATMIL assigns a probability between 0.93
and 0.99 to the same lesions and detects each of them in all five fold
models, so the recovery is not the effect of one favorable training run.
Measured over the whole case, the 13-voxel lesion in P47\_T1 accounts for
0.07\% of the 18,736 lesion voxels of that case. Omitting it therefore has
almost no effect on a voxel-wise overlap term. Under the component-balanced
term and the lesion-level term of CATMIL, however, the same lesion is one
of 59 lesions and carries a comparable share of the lesion-level signal.
Training uses patches rather than whole volumes, so the exact proportions
differ during optimization, but the imbalance is of the same kind.

\textbf{The difficulty of each recovered lesion has a different source.}

\emph{Row~1: a 5-voxel lesion adjacent to another lesion (P47\_T1).} This
lesion is at the resolution limit of the data. It occupies a single axial
slice, and its local contrast of 2.22 is lower than that of 87\% of the
small test lesions. It also lies 3~mm from a larger neighboring lesion, so
its signal is partly mixed with the neighbor's hyperintensity. The Tversky
and FocalTversky models show a weak response ($p \approx 0.25$) and detect
the lesion in one of five models. This is the regime of lesions of at most
10 voxels, in which baseline recall is only 0.12--0.29
(\cref{fig:recall_vs_size_curve}). CATMIL detects the lesion in every model
with $p = 0.99$. The detection is not a clean delineation, however. In
three of the five CATMIL models, the predicted component extends across the
3~mm gap and joins this lesion to its neighbor, and the lesion-wise DSC is
0.28. At this scale, CATMIL establishes that a lesion is present, but not
where its boundary lies relative to an adjacent lesion.

\emph{Row~2: a clearly visible but isolated 13-voxel lesion (P47\_T1).} The
contrast of this lesion (3.62) is at the median of the small test lesions,
and the focal hyperintensity is visible in the FLAIR crop. Its difficulty
is that it is small and isolated: the nearest other lesion is 19~mm away,
so no neighboring lesion provides spatial context, and the lesion
contributes very few voxels to the training signal. All fifteen baseline
models miss it with $p \leq 0.01$. This lesion shows that the baselines'
omission is a size effect rather than a visibility effect. CATMIL detects
it in every model, but its predicted component has on average only 4
voxels, compared with 13 in the reference (lesion-wise DSC 0.42).

\emph{Row~3: a faint 21-voxel lesion next to a detected large lesion
(P47\_T1).} This lesion's contrast (2.56) is in the lowest quartile of the
small lesions. The same crop contains a large lesion, which all four
losses segment; the baselines' omission is specific to the small, faint
component, which receives a probability of 0.00 from all fifteen baseline
models. This lesion illustrates the gain in the 11--50 voxel range, where
CATMIL raises recall from 0.74--0.82 to 0.89. The high-probability region
of CATMIL lies on the upper part of the reference lesion, and its mean
predicted component has 7.6 voxels: CATMIL again localizes the lesion
correctly but covers only part of it.

\emph{Row~4: a juxtacortical lesion (P26\_T1).} This 37-voxel lesion spans
four slices and has median contrast (3.52), but it lies 7~mm from the brain
surface, directly beneath the cortex, which is itself relatively bright on
FLAIR. The ring around this lesion has a median normalized intensity of
0.48, compared with approximately zero for the other recovered lesions, so
the lesion must be separated from cortex-like signal rather than from dark
white matter. DiceCE and Tversky assign $p \approx 0.3$ and detect the
lesion in one of five models, and FocalTversky assigns $p = 0.04$. CATMIL
detects the lesion in every model with $p = 0.98$. Juxtacortical lesions
are one of the lesion locations used to establish dissemination in space in
the diagnostic criteria for multiple sclerosis~\cite{thompson_diagnosis_2018},
so recovering such lesions is clinically relevant beyond the lesion count.

\textbf{The remaining misses are limited by image evidence rather than by
the loss.} The lesion in row~5 (P52\_T2, 18 voxels) is located in the
posterior fossa and is missed by all twenty models, each of which assigns
it a probability of 0.00. Its contrast of 1.49 is lower than that of 99\%
of the small test lesions, and it is surrounded by relatively bright tissue
(median ring intensity 0.95). Both time points are registered to the same
1~mm template space, and the P52\_T1 reference has no lesion within 5~mm of
this location; the lesion is therefore either new at follow-up or an
annotation that is not supported by the earlier scan. CATMIL increases the
weight of each annotated lesion during training but cannot create image
evidence that the input does not contain. The test set as a whole shows
the same pattern: among the 84 lesions of 5--50 voxels, those that CATMIL
still misses (detected by fewer than three of the five models) have a
median contrast of 1.86, compared with 3.65 for the lesions it detects.

\textbf{The false positive results from the same mechanism as the recovered
lesions.} The component in row~6 (P52\_T1) has only 3 voxels, but all five
CATMIL models predict it with $p = 0.90$, whereas no baseline model
predicts it. The FLAIR crop shows a small focal hyperintensity at this
location. The follow-up scan P52\_T2 contains no annotated lesion within
5.8~mm of this location. The component may be a tiny unannotated
hyperintensity or a non-lesional focus such as a vessel or noise; the
reference annotation cannot distinguish between these possibilities. The
property that recovers the lesions in rows~1--4 is a strong response to
compact focal hyperintensities, independent of their volume; the false
positive is the same response at a location where the reference contains
no lesion. Because this component contains fewer than 5 voxels, it is
removed by the connected-component filter evaluated in \cref{sec:fp_pp}.

\textbf{CATMIL extends detection toward smaller and fainter lesions, and it
mainly corrects lesion omission rather than boundary delineation.} Among
the small lesions of 5--50 voxels, the 15 lesions detected by at least
three CATMIL models but by fewer than three DiceCE models have a median
size of 15 voxels and a median contrast of 2.85, whereas the 58 lesions
detected by both losses have a median size of 35 voxels and a median
contrast of 3.83. CATMIL's predictions for the recovered lesions cover only
part of each lesion (lesion-wise DSC 0.28--0.52), and in row~1 they merge
adjacent lesions. The lesion-level term mainly affects whether a lesion is
detected, whereas boundary accuracy continues to depend on the voxel-wise
overlap terms. Consequently, CATMIL changes the dominant error from a
missed lesion to a detected lesion with inaccurate extent, or to a small
additional component.

\subsection{Analysis of False Positives and Post-Processing}
\label{sec:fp_pp}

As shown in \cref{tab:loss_comparison_small_lesion,tab:paired_differences},
CATMIL achieves the highest small-lesion recall with DSC and HD95
comparable to the baselines, but its lesion F1 is lower than those of the
competing methods. CATMIL achieves a high
lesion-level recall of 0.9085, whereas its lesion-level precision is
0.6738: most true lesions are detected, but additional false-positive
components are also introduced.

At first glance, this behavior may suggest over-segmentation. However,
CATMIL produces the lowest mean false-positive volume among all models while
predicting a considerably larger number of connected components,
approximately 51 on average compared with roughly 30 components in the
ground truth. CATMIL therefore does not primarily produce large,
incorrectly segmented lesion regions; it introduces many small
disconnected components, consistent with increased sensitivity to weak
lesion-like signals. The spatial distribution of these false positives
supports this interpretation: when the distance from each false-positive
voxel to the nearest ground-truth lesion is measured, the near-voxel
fraction is approximately 0.7--0.8 and the median distance is only
1--2~mm. Many of the false positives are thus local fragments near lesion
boundaries or lesion-like regions rather than large, remote false alarms.

Based on this analysis, the predictions of all models are post-processed by
removing predicted connected components smaller than 5 voxels, and the
evaluation metrics are recomputed. This analysis is hereafter referred to
as PP5; its results are reported in
\cref{tab:loss_comparison_small_lesion_pp5}.

\begin{table}[t]
\centering
\caption{Comparison of CATMIL with standard losses on the MSLesSeg test
set after post-processing with a minimum connected-component size of 5
voxels (PP5). Best values are shown in \textbf{bold}; second-best values
are \underline{underlined}.}
\label{tab:loss_comparison_small_lesion_pp5}
\resizebox{\textwidth}{!}{%
\begin{tabular}{lcccc}
\toprule
\textbf{Metric}
& \textbf{CATMIL (ours)}
& \textbf{DiceCE}
& \textbf{Tversky}
& \textbf{FocalTversky} \\
\midrule
Small Lesion Recall $\uparrow$
& \textbf{0.8554 $\pm$ 0.01}
& 0.7853 $\pm$ 0.03
& \underline{0.8223 $\pm$ 0.02}
& 0.8196 $\pm$ 0.02 \\
FN Count $\downarrow$
& \textbf{3.1333 $\pm$ 0.21}
& 5.2333 $\pm$ 0.52
& \underline{4.0667 $\pm$ 0.24}
& 4.0833 $\pm$ 0.24 \\
FN Volume Fraction $\downarrow$
& \textbf{0.0238 $\pm$ 0.004}
& 0.0353 $\pm$ 0.01
& 0.0277 $\pm$ 0.01
& \underline{0.0270 $\pm$ 0.01} \\
\midrule
DSC $\uparrow$
& \textbf{0.7842 $\pm$ 0.01}
& 0.7797 $\pm$ 0.01
& 0.7707 $\pm$ 0.01
& \underline{0.7803 $\pm$ 0.01} \\
HD95 (mm) $\downarrow$
& \textbf{8.1028 $\pm$ 1.76}
& 9.0895 $\pm$ 1.37
& 10.3585 $\pm$ 0.21
& \underline{8.2297 $\pm$ 1.95} \\
Lesion F1 $\uparrow$
& 0.8415 $\pm$ 0.01
& \textbf{0.8586 $\pm$ 0.01}
& 0.8532 $\pm$ 0.01
& \underline{0.8582 $\pm$ 0.005} \\
FP Volume (mm$^3$) $\downarrow$
& \textbf{1515 $\pm$ 244}
& \underline{1815 $\pm$ 654}
& 2618 $\pm$ 558
& 2279 $\pm$ 710 \\
\bottomrule
\end{tabular}%
}
\end{table}

PP5 improves lesion F1 for all four methods, with a substantially larger
effect on CATMIL. CATMIL's lesion F1 increases from 0.7571 to 0.8415, an
improvement of 8.44 percentage points, whereas the other methods improve by
only 1.27--1.53 percentage points. Although CATMIL's small-lesion recall
decreases by 1.76 percentage points, from 0.8730 to 0.8554, it remains the
highest among all methods. CATMIL already has the lowest false-positive
volume before post-processing, and PP5 reduces it only slightly, from 1537
to 1515~mm$^3$, while producing a large improvement in lesion F1. The false
positives introduced by CATMIL therefore consist primarily of small
disconnected components that contribute little to the total false-positive
volume but have a considerable effect on lesion-level evaluation.

After PP5, CATMIL's lesion F1 gap relative to the best-performing method
decreases from 8.84 to 1.71 percentage points. In the paired analysis, the
lesion F1 difference after PP5 is no longer supported against any baseline
($-0.017$ [$-0.040$, +0.001] against DiceCE), whereas the small-lesion
recall gain remains present in all six patients against every baseline
(+0.070 [+0.023, +0.151] against DiceCE; \cref{tab:paired_differences}).
Simple post-processing thus substantially mitigates CATMIL's main
limitation without removing its principal advantage.
\Cref{appendix:lesion-level-analyses} extends this analysis from the single
filter size of PP5 to a range of filter sizes and probability thresholds.

\subsection{Detection at Matched Lesion Precision}
\label{sec:matched_precision}

PP5 compares the losses at one fixed setting, at which CATMIL still has a
lower lesion-wise precision than the baselines. A stricter test of whether
CATMIL detects more lesions, rather than only operating at a more liberal
point, is to compare the losses at the same lesion-wise precision. Each
loss was therefore evaluated over a grid of probability thresholds
($t = 0.1, \dots, 0.9$) and minimum component sizes ($s = 1$ to 170
voxels), with lesion-wise precision, recall, and F1 computed as in
\cref{tab:loss_comparison_small_lesion}
(\cref{fig:operating_grid} in \cref{appendix:lesion-level-analyses}). For
each baseline, the setting of CATMIL with the highest small-lesion recall
whose precision is at least that of the baseline at its default output was
selected on MSLesSeg. The same setting was then applied unchanged to
3D-MR-MS, which was not used for the selection
(\cref{tab:matched_precision}).

\begin{table}[t]
\centering
\caption{Small-lesion recall and lesion-wise F1 of CATMIL at matched
lesion-wise precision. For each baseline, the CATMIL setting $(t, s)$ with
the highest small-lesion recall whose precision on MSLesSeg is at least
that of the baseline at its default output $(0.5, 1)$ was selected on
MSLesSeg and applied unchanged to 3D-MR-MS. Values are differences CATMIL
minus baseline with 95\% patient-level bootstrap CIs; differences whose CI
excludes zero are shown in \textbf{bold}. The MSLesSeg values are
optimistic because the setting is selected on the same data.}
\label{tab:matched_precision}
\resizebox{\textwidth}{!}{%
\begin{tabular}{lcccccc}
\toprule
& & \multicolumn{2}{c}{\textbf{MSLesSeg} (selection)} & \multicolumn{3}{c}{\textbf{3D-MR-MS} (transfer)} \\
\cmidrule(lr){3-4}\cmidrule(lr){5-7}
\textbf{Baseline} & \textbf{CATMIL $(t, s)$} & $\Delta$ Small Recall & $\Delta$ Lesion F1 & $\Delta$ Small Recall & $\Delta$ Precision & $\Delta$ Lesion F1 \\
\midrule
DiceCE & (0.1, 15) & \textbf{+0.048 [+0.013, +0.111]} & \textbf{+0.021 [+0.005, +0.041]} & +0.040 [$-$0.002, +0.087] & $-$0.050 [$-$0.099, +0.001] & $-$0.008 [$-$0.036, +0.023] \\
Tversky & (0.8, 4) & +0.016 [$-$0.013, +0.060] & +0.007 [$-$0.017, +0.028] & $-$0.010 [$-$0.049, +0.024] & +0.030 [$-$0.040, +0.094] & +0.008 [$-$0.020, +0.037] \\
FocalTversky & (0.2, 10) & +0.016 [$-$0.012, +0.057] & +0.007 [$-$0.008, +0.023] & +0.006 [$-$0.036, +0.043] & +0.006 [$-$0.050, +0.067] & +0.007 [$-$0.022, +0.040] \\
\bottomrule
\end{tabular}%
}
\end{table}

\textbf{At the precision of DiceCE, CATMIL detects more small lesions and
has a higher lesion-wise F1.} With $t = 0.1$ and $s = 15$, CATMIL reaches
the precision of DiceCE on MSLesSeg (0.853 against 0.852) with a
small-lesion recall that is 0.048 higher and a lesion-wise F1 that is
0.021 higher, both supported in all six patients. The lower lesion F1 of
CATMIL at its default output is therefore a consequence of its operating
point, not of a weaker detector. Transferred unchanged to 3D-MR-MS, the
same setting keeps a recall gain of similar size (+0.040), but its CI
narrowly includes zero, and the precision is 0.050 lower, so the matching
is not preserved exactly across datasets.

\textbf{At matched precision, CATMIL does not detect more lesions than the
Tversky-based losses.} Against Tversky and FocalTversky, the small-lesion
recall gain at matched precision shrinks to +0.016 on MSLesSeg and to
between $-0.010$ and +0.006 on 3D-MR-MS, with CIs that include zero. Part
of the advantage of CATMIL over these losses in
\cref{tab:loss_comparison_small_lesion} therefore comes from its more
sensitive operating point, which the Tversky-based losses can partly reach
by lowering their threshold or filter size. The same holds when every loss
is tuned to its own maximum lesion-wise F1: the best F1 on MSLesSeg is
0.876 for CATMIL and 0.865--0.872 for the baselines, and none of the
differences is supported. The advantage of CATMIL is therefore the
detection gain at the default output and against the standard DiceCE loss,
together with the ability to trade it for precision after training without
retraining the model.

\subsection{Detection of Lesions of at Least 3~mm}
\label{sec:clinical_size}

Current recommendations for MRI in multiple sclerosis state that lesions
should generally be at least 3~mm in longest
diameter~\cite{barkhof_2024_2025,filippi_mri_2016}, and the 2024 revision
of the McDonald criteria introduces the central vein sign as a diagnostic
marker that is assessed only in lesions of at least this
size~\cite{montalban_diagnosis_2025,barkhof_2024_2025}. The small-lesion
threshold of 150 voxels used above is a technical definition, so
\cref{tab:clinical_size} also reports recall for lesions above and below
the 3~mm size, measured as described in the Evaluation Metrics.

\begin{table}[t]
\centering
\caption{Lesion recall by lesion size relative to the 3~mm longest
diameter used in clinical MRI reading. $n$: number of reference lesions.
Recall is the fraction of lesions detected, averaged over the five fold
models. $\Delta$: CATMIL minus DiceCE with 95\% patient-level bootstrap CI;
differences whose CI excludes zero are shown in \textbf{bold}.}
\label{tab:clinical_size}
\resizebox{\textwidth}{!}{%
\begin{tabular}{llrccccl}
\toprule
\textbf{Dataset} & \textbf{Lesions} & $n$ & \textbf{CATMIL} & \textbf{DiceCE} & \textbf{Tversky} & \textbf{FocalTversky} & $\Delta$ vs.\ DiceCE \\
\midrule
MSLesSeg & $\geq$3~mm & 345 & 0.944 & 0.870 & 0.907 & 0.907 & \textbf{+0.074 [+0.032, +0.129]} \\
         & $\geq$3~mm and $\leq$150 voxels & 233 & 0.924 & 0.827 & 0.874 & 0.875 & \textbf{+0.098 [+0.040, +0.171]} \\
         & $<$3~mm & 16 & 0.275 & 0.087 & 0.175 & 0.225 & \textbf{+0.188 [+0.044, +0.300]} \\
\midrule
3D-MR-MS & $\geq$3~mm & 269 & 0.894 & 0.830 & 0.837 & 0.833 & \textbf{+0.065 [+0.015, +0.125]} \\
         & $\geq$3~mm and $\leq$150 voxels & 156 & 0.855 & 0.753 & 0.768 & 0.756 & \textbf{+0.103 [+0.050, +0.203]} \\
         & $<$3~mm & 120 & 0.393 & 0.282 & 0.300 & 0.308 & \textbf{+0.112 [+0.010, +0.167]} \\
\bottomrule
\end{tabular}%
}
\end{table}

\textbf{The detection gain is not limited to lesions below the clinical
size.} For lesions of at least 3~mm, CATMIL detects 94.4\% of the
reference lesions on MSLesSeg, against 87.0\% for DiceCE and 90.7\% for
both Tversky-based losses, and 89.4\% on 3D-MR-MS, against 83.0--83.7\%.
The difference from DiceCE is supported on both datasets. The largest gain
is for small lesions that are still above the clinical size (at least
3~mm and at most 150 voxels), for which CATMIL raises recall by about 10
percentage points on both datasets: these are the lesions that a reader
would count and that remain eligible for central vein sign assessment, but
that a voxel-wise objective tends to miss. Lesions below 3~mm are rare on
MSLesSeg (16 of 361) but common on 3D-MR-MS (120 of 389), which accounts
for part of the lower overall recall on that dataset; recall below 3~mm
remains low for every loss (at most 0.39). These lesions are not
necessarily irrelevant, because lesions smaller than 3~mm in the corpus
callosum or the infratentorial region can still be considered
demyelinating~\cite{barkhof_2024_2025}; lesion location was not analyzed
in this study.

\subsection{Evaluation on a Second Dataset: 3D-MR-MS}

\begin{table}[t]
\centering
\caption{Comparison of CATMIL with standard losses on the 3D-MR-MS test
set. Best values are shown in \textbf{bold}; second-best values are
\underline{underlined}.}
\label{tab:external_eval_3dmrms}
\begin{tabular}{lcccc}
\toprule
\textbf{Metric} & \textbf{CATMIL (ours)} & \textbf{DiceCE} & \textbf{Tversky} & \textbf{FocalTversky} \\
\midrule
Small Lesion Recall $\uparrow$
& \textbf{0.618 $\pm$ 0.03}
& 0.537 $\pm$ 0.01
& \underline{0.564 $\pm$ 0.02}
& 0.563 $\pm$ 0.02 \\
FN Count $\downarrow$
& \textbf{16.9 $\pm$ 0.64}
& 22.0 $\pm$ 0.94
& \underline{21.3 $\pm$ 1.15}
& \underline{21.3 $\pm$ 1.20} \\
HD95 (mm) $\downarrow$
& \textbf{11.55 $\pm$ 1.27}
& \underline{11.94 $\pm$ 1.79}
& 13.38 $\pm$ 3.65
& 12.27 $\pm$ 3.10 \\
FP Volume (mm$^3$) $\downarrow$
& \textbf{1232 $\pm$ 165}
& \underline{1234 $\pm$ 213}
& 1994 $\pm$ 284
& 1958 $\pm$ 269 \\
\midrule
DSC $\uparrow$
& 0.718 $\pm$ 0.01
& \underline{0.730 $\pm$ 0.02}
& 0.726 $\pm$ 0.01
& \textbf{0.734 $\pm$ 0.01} \\
Lesion F1 $\uparrow$
& 0.471 $\pm$ 0.04
& \textbf{0.668 $\pm$ 0.02}
& \underline{0.640 $\pm$ 0.02}
& 0.637 $\pm$ 0.03 \\
FN Volume Fraction $\downarrow$
& 0.085 $\pm$ 0.03
& \textbf{0.068 $\pm$ 0.01}
& 0.082 $\pm$ 0.03
& \underline{0.074 $\pm$ 0.03} \\
\bottomrule
\end{tabular}
\end{table}

The results in \cref{tab:external_eval_3dmrms} show that CATMIL
consistently provides the strongest performance on the metrics most
directly associated with detecting small lesions. CATMIL achieves the
highest small-lesion recall ($0.618 \pm 0.03$), substantially exceeding
DiceCE ($0.537 \pm 0.01$) and both Tversky-based losses ($\approx0.564$).
It also produces the lowest false-negative lesion count
($16.9 \pm 0.64$), reducing the number of completely missed lesions by
approximately 23\% compared with DiceCE ($22.0 \pm 0.94$). Both gains are
present in five of the six test patients against every baseline, with CIs
that exclude zero (\cref{tab:paired_differences}). The false-positive
volume of CATMIL ($1232 \pm 165$~mm$^3$) is similar to that of DiceCE
($1234 \pm 213$~mm$^3$) and lower than that of both Tversky-based losses in
all six patients, so the improved sensitivity is not obtained by
over-segmenting large regions. CATMIL has the lowest mean HD95
($11.55 \pm 1.27$~mm), but none of the HD95 differences is supported.

The cost of this sensitivity is larger on this dataset than on MSLesSeg.
CATMIL has the lowest lesion-wise F1 ($0.471 \pm 0.04$, against
0.637--0.668 for the baselines), lower in every test patient, and a
slightly lower DSC. The DSC difference from DiceCE is small
($-0.012$ [$-0.032$, +0.010]) but has the same sign in all six patients,
so the DSC on this dataset should be regarded as slightly reduced rather
than preserved. Lesion-wise F1
falls when the prediction contains components that overlap no reference
lesion, and without filtering CATMIL produces about 150 such components per
case, against 14--19 for the baselines
(\cref{appendix:lesion-level-analyses}). These components are small, so
they add little false-positive volume, but they lower lesion-wise
precision substantially. CATMIL also has the highest FN Volume Fraction
($0.085 \pm 0.03$), although it misses the fewest lesions; the lesions it
still misses therefore account for a larger share of the lesion volume.

The lesion-level analyses of the previous subsections were repeated on
this dataset (\cref{appendix:lesion-level-analyses}). They show the same
pattern: CATMIL detects about 33 lesions per fold model that DiceCE
misses, against 2.4 in the opposite direction, and almost all of these
recovered lesions are at most 50~mm$^3$. On both datasets, CATMIL improves
the detection of small and sparse lesions without increasing
false-positive volume, and on both the gain is accompanied by additional
small false-positive components. On 3D-MR-MS this cost is large enough
that a component-size filter larger than PP5 is needed to control it. The
lesion-level results on this dataset should also be read with its lesion
distribution in mind: one of the six test patients contributes 249 of the
389 reference lesions, so pooled lesion counts are dominated by this
patient, whereas the patient-level analysis gives each patient equal
weight.
